% =============================================================================
% Section 6: Phase 2 — Supervised Disentangled Projection
% =============================================================================
\section{Phase 2: Supervised Disentangled Projection}
\label{sec:sdp}

% ---------------------------------------------------------------------------
% 6.1 Motivation
% ---------------------------------------------------------------------------
\subsection{Motivation}
\label{sec:sdp:motivation}

The encoder features $\mathbf{h} \in \R^{768}$ are high-dimensional and
unstructured: semantic factors (volume, location, shape) are entangled across
all 768~dimensions. Downstream GP-based growth prediction requires features
that are (i)~\emph{semantically interpretable}: tumour volume should be
decodable from a known subset of latent dimensions; (ii)~\emph{disentangled}:
volume, location, and shape should occupy independent subspaces;
(iii)~\emph{low-dimensional}: the GP covariance matrix scales as $O(D^2)$ in
the output dimensionality $D$.

The Supervised Disentangled Projection (SDP) addresses all three requirements
through a lightweight MLP trained on frozen encoder features with a composite
loss combining semantic supervision, collapse prevention, and independence
regularisation.

\subsubsection{Why Not a VAE?}

Early experiments with a Semi-Supervised VAE (SemiVAE) were abandoned due to
three failure modes: (i)~posterior collapse, where latent dimensions become
uninformative despite extensive hyperparameter tuning (7-phase curriculum
across 600~epochs); (ii)~KL distortion, where the $D_{\KL}(q \| p)$ penalty
actively minimises the mutual information $I(x; z)$ between input and
representation, destroying the rich foundation model features; and
(iii)~reconstruction burden, where the decoder wastes capacity on
pixel-level fidelity irrelevant to downstream growth prediction.

The SDP replaces the generative objective with a discriminative one that
directly optimises for the properties needed by the GP temporal model,
following the theoretical mandate of Locatello et al.\
\cite{locatello2019challenging}: supervised disentanglement resolves the
impossibility of unsupervised approaches.

% ---------------------------------------------------------------------------
% 6.2 Architecture
% ---------------------------------------------------------------------------
\subsection{Architecture}
\label{sec:sdp:architecture}

The SDP consists of a projection network $g_\theta$ and per-partition semantic
heads $\{\pi_p\}$:
\begin{equation}
  g_\theta: \R^{768} \to \R^{128}, \quad
  g_\theta(\mathbf{h}) = \text{SN}(\mathbf{W}_2)\,
    \sigma\bigl(\text{SN}(\mathbf{W}_1)\,
    \text{LN}(\mathbf{h}) + \mathbf{b}_1\bigr) + \mathbf{b}_2,
  \label{eq:sdp_forward}
\end{equation}
where $\text{LN}$ is LayerNorm, $\sigma$ is GELU with dropout $p=0.1$, and
$\text{SN}(\cdot)$ denotes spectral normalisation applied to \emph{both}
linear layers (\cref{sec:bg:disentanglement}). The architecture is:

\begin{equation}
  \underbrace{\mathbf{h} \in \R^{768}}_{\text{frozen encoder}}
  \xrightarrow{\text{LN}}
  \xrightarrow{\text{SN-Linear}(768, 512)}
  \xrightarrow{\text{GELU + Drop}}
  \xrightarrow{\text{SN-Linear}(512, 128)}
  \underbrace{\mathbf{z} \in \R^{128}}_{\text{structured latent}}
  \label{eq:sdp_pipeline}
\end{equation}

The total parameter count is approximately 459K
($768 \times 512 + 512 + 512 \times 128 + 128$).

% ---------------------------------------------------------------------------
% 6.3 Partition Structure
% ---------------------------------------------------------------------------
\subsection{Latent Partition Structure}
\label{sec:sdp:partitions}

The 128-dimensional latent vector is partitioned into semantically supervised
and unsupervised subspaces:

\begin{table}[H]
  \centering
  \caption{Original SDP partition structure ($3+1$ partitions). See
    \cref{sec:eval:revision} for the R1-revised $1+1$ structure.}
  \label{tab:partitions}
  \begin{tabular}{lcccc}
    \toprule
    \textbf{Partition} & \textbf{Dims} & \textbf{Indices} &
      \textbf{Target dim} & \textbf{Semantic head} \\
    \midrule
    $\Vol$ (volume)   & 24 & $[0, 24)$    & 4 & $\pi_{\text{vol}}: \R^{24} \to \R^4$ \\
    $\Loc$ (location) & 8  & $[24, 32)$   & 3 & $\pi_{\text{loc}}: \R^{8} \to \R^3$ \\
    $\Shape$ (shape)  & 12 & $[32, 44)$   & 6 & $\pi_{\text{shape}}: \R^{12} \to \R^6$ \\
    $\Res$ (residual) & 84 & $[44, 128)$  & --- & None (unsupervised) \\
    \midrule
    \textbf{Total}    & 128 & $[0, 128)$ & 13 & --- \\
    \bottomrule
  \end{tabular}
\end{table}

Each supervised partition $p$ has a linear semantic head
$\pi_p: \R^{d_p} \to \R^{k_p}$ that predicts ground-truth semantic features
extracted from segmentation masks. The residual partition captures
information not explained by the three supervised factors.

\begin{remark}[R1 Revision]
  Following empirical results (shape $R^2 \leq 0.11$, temporally static
  meningioma location), the partition structure is revised to
  $\mathbf{z} = [\Vol \in \R^{32} \mid \Res \in \R^{96}]$ with a single
  volume head $\pi_{\text{vol}}: \R^{32} \to \R^1$ predicting
  $\log(V_{\text{WT}} + 1)$. Location is relocated to a static GP covariate
  (\cref{sec:eval:revision}).
\end{remark}

% ---------------------------------------------------------------------------
% 6.4 Loss Function
% ---------------------------------------------------------------------------
\subsection{Composite Loss Function}
\label{sec:sdp:loss}

The SDP training objective combines four terms:
\begin{equation}
  \mathcal{L}_{\SDP} =
    \underbrace{\sum_{p \in \mathcal{P}} \lambda_p\,
      \|\hat{\mathbf{y}}_p - \mathbf{y}_p\|_2^2
    }_{\text{Semantic regression}}
  + \underbrace{\lambda_{\text{var}}\,
      \mathcal{L}_{\text{var}}(\mathbf{z})
    }_{\text{Collapse prevention}}
  + \underbrace{\lambda_{\text{cov}}\,
      \mathcal{L}_{\text{cov}}(\mathbf{z})
    }_{\text{Linear independence}}
  + \underbrace{\lambda_{\text{dCor}}\!
      \sum_{\substack{i,j \in \mathcal{P} \\ i < j}}
      \dCor(\mathbf{z}_i, \mathbf{z}_j)
    }_{\text{Nonlinear independence}},
  \label{eq:sdp_loss}
\end{equation}
where $\mathcal{P}$ is the set of supervised partitions,
$\hat{\mathbf{y}}_p = \pi_p(\mathbf{z}_p)$, and each loss term serves a
distinct purpose.

\subsubsection{Semantic Regression}

For each supervised partition, MSE loss anchors the partition's information
content to a known semantic factor:
\begin{equation}
  \mathcal{L}_p^{\text{sem}} = \frac{1}{k_p}\,
    \|\pi_p(\mathbf{z}_p) - \mathbf{y}_p\|_2^2.
  \label{eq:sem_loss}
\end{equation}
Targets $\mathbf{y}_p$ are normalised to zero mean and unit variance using
statistics computed on the training pool only (\cref{sec:sdp:normalisation}).

\subsubsection{Variance Hinge Loss (Collapse Prevention)}

The VICReg variance term (\cref{eq:vicreg_var}) is applied to all 128
dimensions of $\mathbf{z}$, ensuring no dimension collapses to a constant.
The threshold $\gamma = 1.0$ targets unit standard deviation per dimension.

\subsubsection{Covariance Penalty (Linear Independence)}

Cross-partition covariance is penalised:
\begin{equation}
  \mathcal{L}_{\text{cov}} = \frac{1}{d^2}
    \sum_{\substack{i,j: \\ \text{part}(i) \neq \text{part}(j)}} C_{ij}^2,
  \label{eq:cov_loss}
\end{equation}
where $\mathbf{C}$ is the batch covariance matrix and $\text{part}(i)$
denotes the partition containing dimension $i$. Only cross-partition
off-diagonal elements are penalised; within-partition correlations are
permitted.

\subsubsection{Distance Correlation (Nonlinear Independence)}

For each pair of supervised partitions $(i, j)$, the distance correlation
$\dCor(\mathbf{z}_i, \mathbf{z}_j)$ (\cref{eq:dcor}) captures nonlinear
dependencies not detected by the covariance penalty. The empirical distance
correlation is computed from $O(N^2)$ pairwise distance matrices, which for
$N = 800$ (full-batch training) requires approximately 30~MB of GPU memory
per partition pair.

% ---------------------------------------------------------------------------
% 6.5 Curriculum Training
% ---------------------------------------------------------------------------
\subsection{Curriculum Training Schedule}
\label{sec:sdp:curriculum}

Following the capacity-annealing principle from the disentanglement literature
\cite{burgess2018understanding}, losses are activated in a staged curriculum:

\begin{table}[H]
  \centering
  \caption{SDP curriculum training schedule (4 phases).}
  \label{tab:curriculum}
  \begin{tabular}{clp{7cm}}
    \toprule
    \textbf{Phase} & \textbf{Epochs} & \textbf{Active losses} \\
    \midrule
    Warm-up     & 0--9    & $\mathcal{L}_{\text{var}}$ only \\
    Semantic    & 10--39  & $+ \mathcal{L}_{\text{vol}}, \mathcal{L}_{\text{loc}}, \mathcal{L}_{\text{shape}}$ \\
    Independence & 40--59 & $+ \mathcal{L}_{\text{cov}}, \mathcal{L}_{\text{dCor}}$ \\
    Full         & 60--100 & All losses at full strength \\
    \bottomrule
  \end{tabular}
\end{table}

The rationale is that semantic regression must first establish the partition
structure before independence regularisation is imposed. Premature
independence penalties can collapse dimensions before they have acquired
semantic content.

% ---------------------------------------------------------------------------
% 6.6 Hyperparameters
% ---------------------------------------------------------------------------
\subsection{Hyperparameters}
\label{sec:sdp:hyperparams}

\begin{table}[H]
  \centering
  \caption{SDP training hyperparameters.}
  \label{tab:sdp_hyperparams}
  \begin{tabular}{lcp{6cm}}
    \toprule
    \textbf{Parameter} & \textbf{Value} & \textbf{Rationale} \\
    \midrule
    $\lambda_{\text{vol}}$ & 20.0 & Highest weight; volume is primary target \\
    $\lambda_{\text{loc}}$ & 12.0 & Moderate weight \\
    $\lambda_{\text{shape}}$ & 15.0 & Higher than location (harder to learn) \\
    $\lambda_{\text{cov}}$ & 5.0 & Cross-partition linear independence \\
    $\lambda_{\text{var}}$ & 5.0 & Collapse prevention \\
    $\lambda_{\text{dCor}}$ & 2.0 & Nonlinear independence \\
    Learning rate & $10^{-3}$ & Standard for small MLPs \\
    Optimizer & AdamW & Decoupled weight decay \\
    Weight decay & 0.01 & Regularisation \\
    Epochs & 100 & Sufficient for convergence \\
    Batch size & 800 (full-batch) & Exact covariance/dCor estimates \\
    Scheduler & Cosine, 5-epoch warmup & Smooth convergence \\
    \bottomrule
  \end{tabular}
\end{table}

The full-batch training is computationally trivial since the SDP operates on
precomputed 768-dimensional feature vectors, not 3D volumes. Full-batch
provides exact (not estimated) covariance and distance correlation statistics.

% ---------------------------------------------------------------------------
% 6.7 Normalisation Scope
% ---------------------------------------------------------------------------
\subsection{Normalisation Scope}
\label{sec:sdp:normalisation}

Target normalisation statistics ($\mu$, $\sigma$) are computed on the training
pool (800~subjects) only and applied without recomputation to validation
(100), test (150), and the Andalusian cohort. This prevents information
leakage and ensures consistent scaling across all downstream uses of the
semantic heads.

% ---------------------------------------------------------------------------
% 6.8 Lipschitz Analysis
% ---------------------------------------------------------------------------
\subsection{Lipschitz Continuity Analysis}
\label{sec:sdp:lipschitz}

With spectral normalisation on both linear layers and GELU activation
(Lipschitz constant $\approx 1.13$), the global Lipschitz constant of the SDP
is bounded by:
\begin{equation}
  \operatorname{Lip}(g_\theta) \leq
    \underbrace{1}_{\text{LN}} \cdot
    \underbrace{1}_{\text{SN}(\mathbf{W}_1)} \cdot
    \underbrace{1.13}_{\text{GELU}} \cdot
    \underbrace{1}_{\text{SN}(\mathbf{W}_2)} \approx 1.13.
  \label{eq:lipschitz}
\end{equation}
This ensures that the projection is a near-contraction:
$\|\mathbf{z}_1 - \mathbf{z}_2\| \leq 1.13\,\|\mathbf{h}_1 - \mathbf{h}_2\|$.
Encoder feature noise is not amplified through the projection, yielding
well-conditioned inputs for downstream GP models.

% ---------------------------------------------------------------------------
% 6.9 Quality Targets
% ---------------------------------------------------------------------------
\subsection{Quality Targets}
\label{sec:sdp:targets}

\begin{table}[H]
  \centering
  \caption{SDP quality targets (validation set).}
  \label{tab:sdp_targets}
  \begin{tabular}{lccc}
    \toprule
    \textbf{Metric} & \textbf{Minimum} & \textbf{Target} & \textbf{Blocking?} \\
    \midrule
    Volume $R^2$ & 0.80 & 0.90 & Yes \\
    Location $R^2$ & 0.85 & 0.95 & Yes \\
    Shape $R^2$ & 0.30 & 0.40 & Yes \\
    Max cross-partition correlation & $< 0.30$ & $< 0.20$ & Yes \\
    $\dCor(\text{vol}, \text{loc})$ & $< 0.20$ & $< 0.10$ & Yes \\
    Dims with variance $> 0.5$ & $\geq 85\%$ & $\geq 95\%$ & Yes \\
    \bottomrule
  \end{tabular}
\end{table}

\begin{remark}
  The shape $R^2$ target of 0.30 was not achieved in any experimental
  condition ($R^2 \leq 0.11$), motivating the R1 revision to drop the shape
  partition entirely (\cref{sec:eval:revision}).
\end{remark}
